% Source-to-polynomial conversion for the preferred shared-point systems.
\section{The compiled shared-point polynomials}
\label{sec:compiled-shared-polynomials}

We apply the explicit finite-algebra construction to the preferred
sources \texttt{source-43-conic.json.gz} and
\texttt{source-45-conic-semilocal.json.gz}.  This section proves the
algebraic conversion performed by \texttt{compile\_shared\_source.py}.
The arithmetic argument identifying the source projection with
$\mathbb Z$ is a separate input; it uses the pinned manuscript,
Sun's stated theorem, the scalar normalization of
Section~\ref{subsec:scalar-preparation-zero}, and the degree
refinements of Section~\ref{subsec:shared-degree-refinements}.
The compiler preserves the rational projection of each supplied
source, including every exceptional quadratic case.

\subsection{Retained coordinates and rational parametrization}

The finite-anchor source retains 42 rational coordinates: the
input, eighteen abscissa coordinates for nine points over
$\mathbb Q(\sqrt2)$, two Lucas parameters, one common slope,
eighteen parameters for the nine coordinate poles, and two
parameters for the scalar pole.  Its remaining finite choices
consist of nine point roots, one Lucas root, and ten second
pole roots.  There are twenty formal algebra generators.
The semilocal source retains two additional anchor coordinates
and has one additional finite anchor root.

The first pole roots have already been eliminated by
Lemma~\ref{lem:pole-conic-parametrization}.  More explicitly,
for each nonzero $B$ the conic
\[
                     p^2=B(1-e)(1+3e)
\]
is parametrized using a retained coordinate $t$ and a fixed
choice $\sigma\in\{0,1\}$:
\[
 e=\frac{t^2+(4\sigma-1)B}{t^2+3B},\qquad
 p=-\frac{4Bt(1-\sigma)}{t^2+3B},\qquad t^2+3B\ne0.
\]
The choice zero covers all points with $e\ne1$: its inverse is
$t=p/(e-1)$, and the conic equation gives
$t^2+3B=-4B/(e-1)\ne0$.  The choice one covers $e=1,p=0$;
$t=0$ is permitted because $B\ne0$.  The point $e=-1/3,p=0$
is covered by the choice zero with $t=0$.  No nonzero guard on
$p$ or $t$ is imposed.  Thus this operation preserves the rational
input projection while removing one finite quadratic root in
each of the ten pole blocks.

The fixed nonzero square height scales of
Lemma~\ref{lem:fixed-height-scales} are
\[
                       x_0^2,\qquad (1+x_1^2)^2.
\]
They eliminate the eight component-zero finite choices and their
closed conditions.  Consequently each preferred source has
$16384\cdot2^{10}=2^{24}$ original finite parameter choices:
the scalar correction and the ten conic charts.  These are
fixed product indices, not quantified coordinates.

\subsection{Integral triangular monic equations}

Write $a=a_0+a_1\sqrt2$ and
$a^3-(11-6\sqrt2)a=A+B\sqrt2$.  In the fixed point chart, a
halfpoint ordinate is reconstructed as
\[
 b=C+\frac{B}{2C}\sqrt2,\qquad C\ne0,
 \qquad C^4-AC^2+\frac{B^2}{2}=0.
\]
The source variable $Y=2C$ satisfies the integral monic equation
\[
                         Y^4-4AY^2+8B^2=0.
\]
Its derivative $4Y^3-8AY$ is an explicit source guard.  Hence
all retained point roots are simple, including after base change
to an algebraic closure.  Equivalently, the discriminant of the
quadratic in $Y^2$ is $16(A^2-2B^2)$, and the source guards
imply both $Y\ne0$ and $A^2-2B^2\ne0$.

The Lucas equation has the form
\[
                         \tau^2-a\tau+b=0,
\]
where $a,b$ are integer polynomials in retained parameters.
Set $Z=2\tau-a$ and $\Delta=a^2-4b$.  The equation becomes
$Z^2-\Delta=0$, with inverse $\tau=(Z+a)/2$.  This coordinate
change is invertible over $\mathbb Q$, including when
$\Delta=0$.  A repeated Lucas root is therefore not discarded.

Every other finite equation is already monic and quadratic with
no linear term.  Earlier substitutions can introduce powers of
two in its coefficient.  Suppose its current form is
\[
                          y^2+2^{-k}N=0,
\]
where $N$ is an integer polynomial in retained coordinates and
earlier formal generators.  With $s=\lceil k/2\rceil$ and
$Z=2^s y$, it becomes
\[
                          Z^2+2^{2s-k}N=0.
\]
The exponent $2s-k$ is zero or one.  These changes preserve
triangularity and are invertible over $\mathbb Q$.  The producer
records each exact numerator and power of two.  The independent
checker verifies the identities by induction through the
coefficient DAG, without expanding large polynomials.

\subsection{Quadratic exceptions and geometric simplicity}

For every resulting equation $Z_i^2+c_i=0$, take a fixed binary
choice between
\[
 \begin{array}{c|c|c}
 \text{choice}&\text{monic algebra relation}&\text{additional condition}\\ \hline
 0&Z_i=0&c_i=0,\\
 1&Z_i^2+c_i=0&Z_i\ne0.
 \end{array}
\]
The condition $c_i=0$ is imposed as a closed residual, not as a
quotient relation, so it does not destroy freeness of the
presented algebra.  The second case contributes $Z_i$ to the
guard.  These cases cover precisely all rational roots.  In the
zero case the diagonal derivative is one; in the nonzero case it
is $2Z_i$, nonzero on the guarded locus.

The compiler forms the original guard from its distinct
structural factors, discarding repeated factors and nonzero
integer constants.  The source rational-square clearings $Z=Dq$,
$Z^2=ND$ for $q^2=N/D$ are checked exactly, including the
guard on $D$.  The closed-condition and indexed-anchor argument
denominators are also checked against the recorded guards.
After the substitutions, the
coefficient numerators are integer polynomials and the discarded
numerical denominators are powers of two.

The full guard is the product of the original factors, the
nonzero-case quadratic generators, and the indexed anchor when
present.  Each point derivative is among the original factors.
Thus every diagonal derivative of the triangular monic
presentation is nonzero on the guarded open subset.  At every
geometric point of that subset the fiber has local length one.
This is the exact hypothesis of
Theorem~\ref{thm:explicit-guarded-finite-algebra}; generic
separability alone is not being used.

There are eleven quadratic choices for the finite-anchor source
and twelve for the semilocal source.  A mask with $k$ nonzero
choices has free rank $4^9 2^k$, with respective maxima $2^{29}$
and $2^{30}$.  All source parameters, quadratic masks, separating
integers, and degree cases in the theorem are literal finite
product indices.  Its formal algebra generators are eliminated
by the determinant and coefficient operations.

\subsection{The fixed elliptic exclusion product}

For the finite-anchor source, the guard also contains a fixed
product excluding small elliptic indices.  The integer
\[
                             M=2^{9566463}
\]
and the fixed point $P$ are part of the specification.  Write the
variable abscissa as $(N_0+N_1\sqrt2)/D$ after exact denominator
conversion.  For each finite $nP$, compute its primitive triple
\[
                        x(nP)=\frac{a_n+b_n\sqrt2}{c_n}.
\]
The exact integer factor is
\[
                  (c_nN_0-a_nD)^2-2(c_nN_1-b_nD)^2.
\]
For an identity multiple the factor is one.  The product runs
from $n=1$ to $M$, and $D$ is already guarded.  Since
$\sqrt2\notin\mathbb Q$, a factor vanishes at a rational
specialization exactly when the abscissa equals $x(nP)$.
The triples are defined by the exact binary multiplication
algorithm of Section~\ref{sec:fixed-anchor-height}.  Thus all
integer coefficients and index bounds are specified, with no
additional quantified variables.  The enormous collection of
individual factors is not materialized.

The semilocal source uses its quadratic anchor instead and has
no indexed elliptic exclusion product.  Its complete coefficient
DAG is materialized.  Both variants use the outer indexed
determinants and finite products of the finite-algebra theorem.

\begin{proposition}\label{prop:shared-source-compilation}
The preferred emitted specifications define integer polynomials
$P_{43}$ and $P_{45}$ in respectively 43 and 45 coordinates,
including the input.  Their rational input projections equal
those of their corresponding supplied sources.  They have 42
and 44 existential witnesses, respectively.  Subject to the
arithmetic characterization of those sources, each is a
Diophantine definition of $\mathbb Z$ in $\mathbb Q$.

With $H(P)$ denoting the largest absolute coefficient, the
verified bounds are
\[
 \begin{array}{c|c|c|c}
 &\text{maximum algebra rank}&\deg P&H(P)\\ \hline
 P_{43}&2^{29}&2^{9566680}&2^{\,2^{86098389}}\\
 P_{45}&2^{30}&2^{243}&2^{\,2^{246}}.
 \end{array}
\]
The degree and height entries are upper bounds, not asserted
exact values.
\end{proposition}
\begin{proof}
The conic parametrization covers every rational source solution,
and all subsequent rational coordinate changes are invertible.
The two quadratic cases cover all roots and impose exactly the
original equations.  Denominator coverage makes the cleared
source equations and guards equivalent to the original ones.
The fixed anchor clearing is equivalent by the norm identity
just proved.  Hence every finite branch has its intended
rational projection.  The diagonal derivative check permits
application of Theorem~\ref{thm:explicit-guarded-finite-algebra}.
Its one added coordinate gives the stated counts, and a finite
product realizes the union of all source choices and masks.

For the bounds, apply the weighted reduction and multiplication
matrix norm estimates of Section~\ref{subsec:finite-algebra-bounds}
with nine quartics and respectively eleven or twelve quadratic
choices.  Proposition~\ref{prop:fixed-exclusion-norm} bounds the
indexed fixed exclusion product.  Every generator weight,
operator norm estimate, determinant and resultant bound, and
finite factor count is recorded in the specification and checked
independently.  They give the displayed rounded estimates.  For
the indexed exclusion product the estimates are stored as exact
powers of two, so the bound computation itself does not require
materializing integers with millions of digits.
\end{proof}

The specifications contain 6224 and 6245 coefficient-DAG nodes.
They are exact compact descriptions using integer polynomial
operations, coefficient extraction, determinants, and explicitly
bounded finite products.  They are not expanded coefficient
tables or circuits listing every gate of rank-sized matrices.
An independent exact interpreter executes the complete theorem
on small algebras, including nilpotents, root collisions, degree
drops, and zero quadratic cases.  The theorem and source conversion
proof justify the unexpanded large instances.

The other emitted source variants are preserved in separate
packages for independent review of the arithmetic replacements.
All use the same compiler and the same proof of algebraic
conversion.  The package index records their individual counts,
bounds, hashes, and passing verification receipts.
